\chapterbackmatter{Weinberg Exercises}

\begin{enumerate}
\WeinbergExercise{1} Consider a non-relativistic particle of mass $m$, moving along the
  $x$-axis in a potential $V(x)=m\omega^2x^2/2$. Use path-integral methods
  to find the probability that if the particle is at $x_1$ at time $t_1$,
  then it is between $x$ and $x+dx$ at time $t$.

\WeinbergExercise{2} Find the wave function in field space of a state consisting of a single
  spinless particle of mass $m\neq0$. Use the result to derive the Feynman
  rules for emission or absorption of such a particle.

\WeinbergExercise{3} Find the wave function in field space of the vacuum in the theory of a
  neutral vector field of mass $m\neq0$. Use the result to derive the form of
  the $i\epsilon$ terms in the propagator of this field.

\WeinbergExercise{4} The Lagrangian density of the free spin $3/2$ Rarita--Schwinger field
  $\psi^\mu$ is
  \[
    \mathcal{L}=-\bar\psi^\mu(\gamma^\nu\partial_\nu+m)\psi_\mu
    -\frac{1}{3}\bar\psi^\mu(\gamma_\mu\partial_\nu
      +\gamma_\nu\partial_\mu)\psi^\nu
    +\frac{1}{3}\bar\psi^\mu\gamma_\mu
      (\gamma^\sigma\partial_\sigma-m)\gamma^\nu\psi_\nu.
  \]
  Use path-integral methods to find the propagator of this field.
\end{enumerate}
